% Seccion: Parámetros termo-fluidodinámicos relevantes
\section{Parámetros termo-fluidodinámicos relevantes}
\label{sec:cap02_parametros}

El número de Reynolds es el parámetro adimensional clave para caracterizar el estado cinemático del flujo en el serpentín.
El número de Nusselt representa la relación entre la transferencia de calor convectiva y la conducción del propio fluido.
La eficiencia térmica instantánea se define como la fracción de radiación solar útil absorbida y transferida al aire.
La exergía perdida cuantifica la degradación del potencial de trabajo útil debido al rozamiento viscoso y los gradientes térmicos infinitos.
El factor de pérdidas global caracteriza la tasa de escape de calor hacia el ambiente a través del tope, lados y fondo.

\begin{cajaadvertencia}{Pérdidas Térmicas y de Exergía}
De acuerdo con la segunda ley de la termodinámica, maximizar la eficiencia térmica requiere minimizar la generación de entropía.
Las pérdidas exergéticas en colectores solares de aire son originadas principalmente por la caída de presión en curvas cerradas.
\end{cajaadvertencia}
